\documentclass[11pt,a4paper]{article}

\usepackage[margin=25mm]{geometry}
\usepackage{amsmath,amssymb}
\usepackage{booktabs}
\usepackage{hyperref}
\usepackage{listings}
\usepackage{xcolor}

\lstset{
  basicstyle=\ttfamily\small,
  breaklines=true,
  columns=fullflexible,
  frame=single,
  showstringspaces=false
}

\title{Refactoring a Java Fourier-Domain Circuit Simulator\\
for Nonlinear Transistor Support}
\author{Architecture and Implementation Guide}
\date{}

\begin{document}
\maketitle

\section{Purpose and scope}

The existing simulator represents signals using selected Fourier
coefficients, models linear components by complex impedances, and solves
a linear system at each frequency.

The proposed extension is:
\begin{enumerate}
  \item Assemble linear components using admittance stamps and modified
        nodal analysis (MNA).
  \item Preserve the existing independent-frequency linear solver.
  \item Introduce nonlinear device models defined by terminal currents,
        terminal charges, and their derivatives.
  \item Solve nonlinear periodic steady states using harmonic balance.
  \item Adapt the active frequency set using omitted-equation residuals
        and estimated voltage errors.
\end{enumerate}

Do not replace nonlinear transistors with fixed small-signal impedances:
that approach cannot generate harmonics or intermodulation.

This guide targets steady-state simulation. General transients,
autonomous oscillators, subharmonic solutions, and chaotic behavior
require additional methods.

\section{Conventions}

Use the Fourier convention
\[
  v(t)=\sum_{k\in\mathcal K}V_k e^{j k\omega_0t}.
\]
For real signals,
\[
  V_{-k}=V_k^*, \qquad V_0\in\mathbb R.
\]

Use current injected into a network node as positive on the
right-hand side. Device terminal currents are positive into the device.

Choose a ground node and omit its voltage unknown. Always retain DC,
even when all independent sources are AC.

Internally, represent the independent Fourier unknowns using real
numbers:
\[
  x=
  [X_0,\Re X_1,\Im X_1,\ldots]^T.
\]
Here $X_k$ contains node voltages and any MNA branch-current unknowns.
Generate negative-frequency coefficients by conjugation.

Do not store both conjugate coefficients as independent unknowns.

\section{Linear component assembly}

\subsection{Two-terminal admittance stamp}

For a component between nodes $a$ and $b$,
\[
  y(\omega)=\frac{1}{Z(\omega)}, \qquad
  I_{ab}=y(V_a-V_b).
\]

Stamp:
\[
\begin{aligned}
 Y_{aa}&\mathrel{+}=y,& Y_{ab}&\mathrel{-}=y,\\
 Y_{ba}&\mathrel{-}=y,& Y_{bb}&\mathrel{+}=y.
\end{aligned}
\]

Skip entries involving the ground row or column.

\begin{center}
\begin{tabular}{lll}
\toprule
Component & Admittance & DC treatment\\
\midrule
Resistor & $1/R$ & Same stamp\\
Capacitor & $j\omega C$ & Open circuit\\
Inductor & $1/(j\omega L)$ & Use MNA\\
General impedance & $1/Z(\omega)$ & Handle singular limits explicitly\\
\bottomrule
\end{tabular}
\end{center}

Do not approximate ideal shorts using arbitrarily large admittances.

\subsection{Independent current source}

A current source $I_s$ directed from $a$ to $b$ contributes
\[
  S_a\mathrel{-}=I_s,\qquad S_b\mathrel{+}=I_s.
\]
For periodic sources, use their Fourier coefficients at each frequency.

\subsection{Modified nodal analysis}

Use MNA for ideal voltage sources, inductors, and other constraints
that require branch-current unknowns.

For a branch from $a$ to $b$, let $b$ denote its incidence vector:
$+1$ at node $a$, $-1$ at node $b$, and zero elsewhere.

An ideal voltage source with unknown current $I_s$ gives
\[
\begin{bmatrix}
Y & b\\
b^T & 0
\end{bmatrix}
\begin{bmatrix}
V\\ I_s
\end{bmatrix}
=
\begin{bmatrix}
S\\ V_s
\end{bmatrix}.
\]

An inductor with current $I_L$ satisfies
\[
 b^T V-j\omega L I_L=0,
\]
and contributes the block
\[
\begin{bmatrix}
Y & b\\
b^T & -j\omega L
\end{bmatrix}.
\]
At DC this enforces a short circuit without division by zero.

Call the complete linear MNA matrix $A_{\mathrm L}(\omega)$.
It need not be symmetric.

\subsection{Voltage-controlled current source}

For an output current from $a$ to $b$,
\[
 I_{ab}=g(V_c-V_d),
\]
stamp
\[
\begin{aligned}
Y_{ac}&\mathrel{+}=g,&Y_{ad}&\mathrel{-}=g,\\
Y_{bc}&\mathrel{-}=g,&Y_{bd}&\mathrel{+}=g.
\end{aligned}
\]

This is the basic building block of small-signal transistor models.

\section{Nonlinear device contract}

For each device, define reduced terminal voltages
\[
 u=B V,
\]
where $B$ extracts voltage differences from global node voltages.

Define reduced conductive currents and charges:
\[
 f(u), \qquad q(u).
\]

The device contributes node currents
\[
 i_{\mathrm{node}}
 =B^T\left(f(u)+\frac{d}{dt}q(u)\right).
\]

The model must also provide
\[
 G(u)=\frac{\partial f}{\partial u},\qquad
 C(u)=\frac{\partial q}{\partial u}.
\]

The mapping and current definitions must be paired correctly.
Reduced currents are not necessarily currents entering the terminals
used to define the voltage differences.

For a first implementation, allow $q=0$ and $C=0$.
Add charge storage only after conductive models are validated.

Charge-based modeling is preferred to inserting an arbitrary
voltage-dependent capacitance: the physical current is $dq/dt$.

\section{Initial transistor model: Ebers--Moll BJT}

For an NPN device, define
\[
 u=
 \begin{bmatrix}
 v_B-v_E\\
 v_B-v_C
 \end{bmatrix}.
\]

Let
\[
 I_F=I_{ES}\left(e^{v_{BE}/V_T}-1\right),\qquad
 I_R=I_{CS}\left(e^{v_{BC}/V_T}-1\right),
\]
where
\[
 V_T=\frac{k_B T}{q_e}.
\]

Terminal currents positive into the transistor are
\[
\begin{aligned}
 i_C&=\alpha_F I_F-I_R,\\
 i_E&=-I_F+\alpha_R I_R,\\
 i_B&=(1-\alpha_F)I_F+(1-\alpha_R)I_R.
\end{aligned}
\]
Verify
\[
 i_B+i_C+i_E=0.
\]

For this voltage mapping, use the reduced current vector
\[
 f(u)=
 \begin{bmatrix}
 -i_E\\ -i_C
 \end{bmatrix}
 =
 \begin{bmatrix}
 I_F-\alpha_R I_R\\
 -\alpha_F I_F+I_R
 \end{bmatrix}.
\]

With
\[
 g_F=\frac{I_{ES}}{V_T}e^{v_{BE}/V_T},\qquad
 g_R=\frac{I_{CS}}{V_T}e^{v_{BC}/V_T},
\]
the derivative is
\[
 G(u)=
 \begin{bmatrix}
 g_F&-\alpha_R g_R\\
 -\alpha_F g_F&g_R
 \end{bmatrix}.
\]

For a physically consistent Ebers--Moll parameterization, enforce
\[
 \alpha_F I_{ES}=\alpha_R I_{CS},
 \qquad 0<\alpha_F,\alpha_R<1.
\]

This initial model omits Early effect, breakdown, and detailed charge
storage. Document these limitations.

Use numerically stable exponential evaluation, including
\texttt{expm1} near zero. Use Newton step limiting or a documented,
derivative-consistent exponential continuation to prevent overflow.
Do not silently clamp currents while retaining inconsistent derivatives.

\section{Harmonic-balance equations}

For each retained harmonic,
\[
 R_k=
 A_{\mathrm L}(k\omega_0)X_k+
 I_{\mathrm{NL},k}(X)-S_k=0.
\]

Nonlinear currents enter the node-equation rows. Their contribution
to auxiliary branch-constraint rows is zero unless a particular device
model explicitly introduces additional equations.

Linear components remain independent between frequencies.
Nonlinear components couple frequencies.

\subsection{Alternating frequency/time evaluation}

For period $T=2\pi/\omega_0$, choose
\[
 t_m=\frac{mT}{M},\qquad m=0,\ldots,M-1.
\]

Reconstruct terminal voltages:
\[
 u_{d,m}=
 \sum_{\ell\in\mathcal K}B_d V_\ell
 e^{j\ell\omega_0t_m}.
\]

Evaluate
\[
 f_{d,m}=f_d(u_{d,m}),\qquad
 q_{d,m}=q_d(u_{d,m}).
\]

Transform back:
\[
 F_{d,k}=\frac1M\sum_m f_{d,m}e^{-jk\omega_0t_m},
 \qquad
 Q_{d,k}=\frac1M\sum_m q_{d,m}e^{-jk\omega_0t_m}.
\]

Then
\[
 \boxed{
 I_{\mathrm{NL},k}
 =\sum_d B_d^T
 \left(F_{d,k}+jk\omega_0Q_{d,k}\right).
 }
\]

For sparse active spectra, zero-fill inactive coefficients on the FFT
grid. Document the FFT library's normalization and index conventions.

A grid capable of reconstructing the voltage spectrum is not necessarily
large enough to resolve the nonlinear current spectrum. Oversample and
increase $M$ until retained coefficients and candidate error indicators
stabilize. Exponential models have no finite exactly alias-free grid.

\subsection{Newton Jacobian}

At every time sample evaluate $G_d(u)$ and $C_d(u)$, then transform
these derivative matrices to obtain coefficients $G_{d,p}$ and
$C_{d,p}$.

For a signed-frequency representation, the node-voltage blocks are
\[
 \boxed{
 J_{k\ell}
 =
 \delta_{k\ell}A_{\mathrm L}(k\omega_0)
 +\sum_d B_d^T
 \left(G_{d,k-\ell}+jk\omega_0C_{d,k-\ell}\right)B_d.
 }
\]

In this formula, embed $B_d$ and the nonlinear stamps into the full MNA
layout as necessary. The charge multiplier is the output frequency
$jk\omega_0$, not the input frequency $j\ell\omega_0$.

Solve
\[
 J\Delta x=-R,
 \qquad x_{\mathrm{new}}=x+\lambda\Delta x,
 \qquad 0<\lambda\leq1.
\]

For a real packed spectrum, include the dependence of negative-frequency
coefficients on conjugate positive-frequency coefficients. Do not simply
discard the negative-frequency Jacobian blocks.

A practical alternative is a Jacobian-vector product:
\begin{enumerate}
 \item Reconstruct the voltage perturbation in time.
 \item Compute $\delta f_m=G_m\delta u_m$ and
       $\delta q_m=C_m\delta u_m$.
 \item Transform back and form
       $\delta F_k+jk\omega_0\delta Q_k$.
 \item Add the linear MNA product.
\end{enumerate}

Initially use a directly assembled Jacobian for small test circuits.
Introduce matrix-free Newton--Krylov methods only after validation.

\subsection{Convergence and robustness}

Use:
\begin{itemize}
 \item A DC operating-point solve before AC continuation.
 \item Source stepping from small AC amplitude to the target amplitude.
 \item Damped Newton steps with a scaled residual merit function.
 \item Separate absolute and relative tolerances for voltage, current,
       and constraint equations.
 \item Both residual and update-size convergence checks.
 \item Iteration limits and explicit failure diagnostics.
\end{itemize}

MNA rows have different units. Never compare all raw residual entries
using one unscaled tolerance.

Detect floating nodes, inconsistent ideal-source constraints, and
singular matrices. Do not silently hide these problems with arbitrary
regularization.

\section{Adaptive frequency selection}

\subsection{Why source-only selection fails}

Nonlinearities generate DC, harmonics, and mixing products. For example,
$v^2$ produces sums and differences of input frequencies.
Transistor exponentials can generate arbitrarily high orders.

Adaptation must measure the equations omitted by the active spectrum,
rather than only inspect source types.

\subsection{Active and candidate sets}

Maintain:
\begin{itemize}
 \item An active set $\mathcal K$ used by the nonlinear solve.
 \item A larger candidate set $\mathcal P$ inspected for spectral error.
 \item A time-sampling grid sufficiently large to inspect $\mathcal P$.
\end{itemize}

Always retain DC. Add and remove conjugate pairs together.

Maintain source coefficients on the candidate set even when they are
not active. This matters for truncated square and sawtooth sources.

\subsection{Omitted-equation residual}

After convergence on $\mathcal K$, evaluate nonlinear currents on the
larger grid. For omitted $p$, set $X_p=0$ and compute
\[
 \boxed{r_p=I_{\mathrm{NL},p}-S_p.}
\]

This residual may include nonzero voltage-source constraint entries
through $S_p$.

Current magnitude alone is not a reliable spectral error measure:
small currents near resonances can cause large voltages.

Estimate the missing solution by
\[
 \Delta\widehat X_p=-D_p^{-1}r_p,
\]
where
\[
 D_p=A_{\mathrm L}(\omega_p)
 +\operatorname{embed}
 \left[
 \sum_d B_d^T(G_{d,0}+j\omega_pC_{d,0})B_d
 \right].
\]

This is an approximate incremental loading calculation. It ignores
cross-frequency corrections.

For node-voltage entries define
\[
 \eta_p=
 \max_n
 \frac{|\Delta\widehat V_{n,p}|}
 {V_{\mathrm{abs}}+V_{\mathrm{rel}}V_{\mathrm{scale},n}}.
\]

Add separate current or observable error indicators if needed.
A voltage-only indicator does not bound every branch-current error.

If the estimator block is singular or ill-conditioned, do not classify
the candidate as harmless. Expand the candidate solve, improve the
coupled estimator, or report an unresolved error.

\subsection{Adaptation loop}

\begin{lstlisting}
active = initialSourceHarmonicsPlusDC()
x = solveDCAndInitialize(active)

repeat:
    x = solveHarmonicBalance(active, x)
    ensureSamplingConvergence(x)

    candidates = buildGuardSpectrum(active, sources)
    residuals = evaluateOmittedEquations(x, candidates)
    indicators = estimateMissingSolutionErrors(residuals)

    additions = selectWorstCandidates(indicators)
    if additions is empty and guardSpectrumIsAdequate:
        return convergedResult

    active = active union additions
    x = warmStartWithPredictedCoefficients(x, additions)

until resourceLimit

return explicitNonconvergenceDiagnostics
\end{lstlisting}

Inspect beyond the immediate active spectrum and enlarge the guard band
until its outer edge is insignificant. A small residual inside a narrow
candidate set does not prove that more distant harmonics are negligible.

For removal, require:
\begin{enumerate}
 \item Small contribution for several adaptation cycles.
 \item A removal threshold lower than the addition threshold.
 \item An omitted-residual test after trial removal.
 \item Re-solution of the reduced system.
\end{enumerate}

Never remove DC or explicitly protected source components.

Declare separate outcomes for nonlinear convergence, spectral
convergence, and sampling convergence.

\section{Frequency representation}

\subsection{Commensurate sources}

For a common fundamental, use integer indices
\[
 \omega_k=k\omega_0.
\]
Do not identify harmonics by floating-point equality.

Sparse active coefficients can coexist with a dense FFT workspace.

\subsection{Incommensurate sources}

For independent fundamentals $\Omega_1,\ldots,\Omega_d$, use integer
tuples:
\[
 \omega_{\mathbf k}=\mathbf k\cdot\boldsymbol\Omega,
 \qquad \mathbf k\in\mathbb Z^d.
\]

Examples:
\[
 (1,0)\mapsto\Omega_1,\quad
 (1,-1)\mapsto\Omega_1-\Omega_2,\quad
 (2,1)\mapsto2\Omega_1+\Omega_2.
\]

Evaluate the quasiperiodic waveform on a multidimensional phase grid.
The charge multiplier becomes $j\omega_{\mathbf k}$.
Use multidimensional FFTs or equivalent Fourier projection.

Reduce rationally dependent fundamentals to an independent basis.
Otherwise different tuples can describe the same physical frequency.

Implement the one-fundamental periodic solver first. Multitone
adaptation is a later milestone because its frequency lattice can grow
rapidly.

\section{Suggested Java architecture}

Adapt names and types to the existing repository. Avoid introducing
duplicate complex-number or sparse-matrix libraries unnecessarily.

\begin{lstlisting}[language=Java]
interface LinearComponent {
    void allocateUnknowns(UnknownAllocator allocator);
    void stamp(
        double angularFrequency,
        ComplexMatrixAccumulator matrix,
        ComplexVectorAccumulator rhs,
        UnknownLayout layout);
}

interface NonlinearDevice {
    TerminalMapping mapping();

    // Currents and charges use the reduced terminal coordinates.
    void evaluate(
        double[] terminalVoltages,
        double[] currents,
        double[] charges,
        double[][] currentJacobian,
        double[][] chargeJacobian);
}

interface PeriodicSource {
    Complex coefficient(int harmonicIndex);
}

interface SpectrumTransform {
    // Enforce documented normalization and conjugate symmetry.
    void toTime(Spectrum spectrum, TimeSamples output);
    void toFrequency(TimeSamples samples, Spectrum output);
}
\end{lstlisting}

Recommended responsibilities:
\begin{description}
 \item[UnknownLayout] Ground handling, node indices, auxiliary currents,
       and real-spectrum packing.
 \item[LinearAssembler] Frequency-dependent MNA stamps and source vectors.
 \item[DeviceEvaluator] Time-domain device evaluation and derivatives.
 \item[HarmonicBalanceResidual] Transform orchestration and residual assembly.
 \item[NewtonSolver] Scaling, damping, convergence, and diagnostics.
 \item[SpectrumController] Candidate generation, error estimation,
       additions, removals, and hysteresis.
 \item[SamplingController] FFT grid selection and aliasing checks.
\end{description}

Cache linear matrices when frequency and component parameters are
unchanged. Reuse work arrays and sparse symbolic factorizations where
supported. Never reuse numerical factorizations after matrix entries
change.

\section{Implementation milestones for Copilot}

\subsection*{Milestone 1: inspect and preserve behavior}
\begin{enumerate}
 \item Inspect the existing graph representation, ground convention,
       Fourier normalization, solver, and component APIs.
 \item Identify tests and establish baseline results.
 \item Propose a minimal migration plan before large edits.
 \item Preserve public APIs where feasible.
\end{enumerate}

\subsection*{Milestone 2: linear stamps and MNA}
\begin{enumerate}
 \item Wrap existing scalar impedances with a two-terminal stamp.
 \item Add current-source and voltage-source stamps.
 \item Add branch-current treatment for inductors, including DC.
 \item Validate existing linear simulations before introducing nonlinear
       frequency coupling.
\end{enumerate}

\subsection*{Milestone 3: DC nonlinear solver}
\begin{enumerate}
 \item Add nonlinear device current and derivative contracts.
 \item Implement a diode and conductive Ebers--Moll NPN.
 \item Implement damped Newton and source stepping.
 \item Validate analytic derivatives against finite differences.
\end{enumerate}

\subsection*{Milestone 4: fixed-spectrum harmonic balance}
\begin{enumerate}
 \item Implement real Fourier packing and transform tests.
 \item Implement time-domain nonlinear evaluation.
 \item Implement residual and coupled Jacobian assembly.
 \item Solve a user-specified harmonic set.
 \item Validate against enlarged fixed-spectrum solutions.
\end{enumerate}

\subsection*{Milestone 5: adaptive spectrum}
\begin{enumerate}
 \item Add guard-spectrum evaluation and omitted residuals.
 \item Add incremental missing-voltage estimation.
 \item Add batch expansion and warm starts.
 \item Add sampling refinement.
 \item Add conservative removal only after addition is reliable.
\end{enumerate}

\subsection*{Milestone 6: charge and advanced models}
\begin{enumerate}
 \item Add documented terminal charge models and derivatives.
 \item Validate the output-frequency charge multiplier.
 \item Add improved BJT or MOSFET models as separate implementations.
 \item Consider multitone and matrix-free solvers after the periodic
       implementation is stable.
\end{enumerate}

Each milestone should compile, include tests, and avoid unrelated
refactoring.

\section{Required verification tests}

\subsection{Linear assembly}
\begin{itemize}
 \item Resistor divider.
 \item RC and RL frequency responses.
 \item Voltage-source polarity and current-source direction.
 \item Inductor at DC.
 \item Controlled-source signs.
 \item Explicit detection of floating or inconsistent circuits.
 \item Agreement with the original solver for supported circuits.
\end{itemize}

\subsection{Device models}
\begin{itemize}
 \item BJT terminal-current conservation.
 \item Finite-difference verification of every derivative entry.
 \item Reverse-biased and forward-biased junction behavior.
 \item Bias points spanning cutoff, forward active operation, and saturation.
 \item Consistency of charge derivatives when charge models are added.
\end{itemize}

\subsection{Fourier and harmonic balance}
\begin{itemize}
 \item FFT round trip, DC normalization, and conjugate symmetry.
 \item For $v=A\cos(\omega_0t)$ and $i=a v^2$, verify
       $I_0=aA^2/2$ and $I_{\pm2}=aA^2/4$.
 \item A linear circuit solved by harmonic balance must reproduce the
       independent-frequency solver.
 \item A small nonlinear perturbation must approach the corresponding
       small-signal result.
 \item A driven nonlinear circuit must generate DC and harmonics.
 \item Compare Jacobian-vector products with finite-difference residual
       perturbations, including charge terms.
 \item Increasing the sampling grid must stabilize retained results.
\end{itemize}

\subsection{Adaptive spectrum}
\begin{itemize}
 \item Begin with DC and the driven fundamental; verify generated
       harmonics are added.
 \item Use a resonant network to verify that a small omitted current with
       a large voltage effect is retained.
 \item Compare adaptive solutions with generously sized fixed-spectrum
       reference solutions.
 \item Verify tighter tolerances reduce observable error.
 \item Verify removal hysteresis prevents repeated add/remove cycles.
 \item Verify resource limits report unresolved spectral error rather
       than claiming convergence.
\end{itemize}

\section{Acceptance criteria}

The refactor is complete when:
\begin{enumerate}
 \item Existing linear results remain correct.
 \item Ideal voltage sources and inductors work at DC using MNA.
 \item Nonlinear transistors generate DC shifts and harmonics.
 \item Newton derivatives pass numerical verification.
 \item Adaptive solutions agree with expanded-spectrum references within
       configured tolerances.
 \item Aliasing checks are independent of active-set convergence.
 \item Unsupported regimes and convergence failures are clearly reported.
\end{enumerate}

\section{Limitations to document}

Harmonic balance solves an assumed periodic or quasiperiodic steady
state. Convergence does not prove dynamic stability or uniqueness.

Autonomous oscillators require an unknown oscillation frequency and a
phase condition. Subharmonic solutions may require a lower fundamental
than the source frequency. General startup behavior and chaotic signals
require time-domain simulation or other specialized methods.

Adaptive frequency indicators are practical error estimates, not a
universal rigorous bound. Validate their effectiveness against larger
spectra and refined sampling grids.

\end{document}